% Hikaru.
\subsection{\expt{a112} and \expt{a113}: Read-only declarations}
Experiment a112 adds \texttt{READ ONLY} to T3, T4 and T6 on top of a111.
Their query bodies and \texttt{SERIALIZABLE} isolation remain unchanged.
Experiment a113 additionally declares them \texttt{DEFERRABLE}. PostgreSQL permits
a serializable, read-only, deferrable transaction to wait for a safe snapshot.
That property applies to the qualifying transaction, not to every transaction
in the distributed workload \citep{pg-set-transaction}.

\begin{table}[htbp]
\caption{Four-client diagnostics: 60-second warmups and measurements.}
\label{tab:a-read-only}
\begin{tabularx}{\linewidth}{@{}llYr@{}}
\toprule
\textbf{Experiment} & \textbf{Job} & \textbf{Change} & \textbf{TPS} \\
\midrule
\expt{a111} & 977617 & Indexes and join-key colocation & 89.917 \\
\expt{a112} & 977618 & a111 plus READ ONLY & 92.383 \\
\expt{a113} & 977619 & a112 plus DEFERRABLE & 90.550 \\
\bottomrule
\end{tabularx}
\end{table}

All three completed both phases without measured failures or retries and passed
validation. Their small differences do not establish a reliable ranking from
single short runs. The table neither isolates a substantial READ ONLY benefit
nor demonstrates a general concurrency fix. We selected a112 for full evaluation;
the deferrable variant did not demonstrate an additional performance benefit.
